\documentclass[journal, twocolumn]{IEEEtran}
\usepackage{cite}
\usepackage{amsmath,amssymb,amsfonts}
\usepackage{algorithmic}
\usepackage{graphicx}
\usepackage{textcomp}
\usepackage{xcolor}
\usepackage{booktabs}
\usepackage{multirow}
\usepackage{hyperref}

\begin{document}

\title{Resolution Preservation over Architectural Complexity in Pancreas Segmentation from CT Images}

\author{Istiaque Ahmed,
        Kursat Komurcu
\thanks{I. Ahmed and K. Komurcu are with the Institute of Computer Science, Vilnius University, Lithuania.}
}

\markboth{IEEE Transactions on Medical Imaging,~Vol.~X, No.~X, Month~202X}%
{Ahmed \MakeLowercase{\textit{et al.}}: Resolution Preservation over Architectural Complexity in Pancreas Segmentation}

\maketitle

\begin{abstract}
The automated segmentation of the pancreas from Computed Tomography (CT) images remains a significant challenge due to the organ's small volume, high anatomical variability, and low contrast with surrounding tissues. Recent trends in medical image segmentation have heavily favored complex architectures, such as Vision Transformers and Attention-based networks, often at the expense of spatial resolution to accommodate GPU memory constraints. In this work, we empirically demonstrate that for small-organ segmentation, preserving native spatial resolution is substantially more critical than architectural complexity. We propose a simple, high-resolution patch-based Convolutional Neural Network (CNN) framework paired with domain-specific Hounsfield Unit (HU) windowing. Our comprehensive ablation studies reveal that downsampling CT slices from $512 \times 512$ to $256 \times 256$ causes a catastrophic collapse in segmentation fidelity, regardless of the global context provided. Our optimized baseline achieves a state-of-the-art Dice similarity coefficient of 0.849 on the NIH Pancreas-CT dataset, outperforming downsampled full-slice models which stagnate at $\sim$0.31 Intersection-over-Union (IoU). These findings challenge the prevailing shift towards heavy transformer models for this task, advocating for data-centric, resolution-preserving strategies.
\end{abstract}

\begin{IEEEkeywords}
Pancreas Segmentation, Computed Tomography, Deep Learning, Resolution Preservation, Convolutional Neural Networks, Vision Transformers.
\end{IEEEkeywords}

\section{Introduction}
\IEEEPARstart{P}{ancreatic} cancer has one of the lowest survival rates among all solid tumors, largely due to late diagnosis. Automated segmentation of the pancreas from abdominal Computed Tomography (CT) scans is a crucial prerequisite for computer-aided diagnosis and surgical planning. However, it is notoriously difficult. The pancreas occupies a minute fraction of the abdominal volume, exhibits high shape variability, and shares similar Hounsfield Unit (HU) intensities with neighboring organs like the stomach and duodenum.

To address these challenges, deep learning approaches, particularly the U-Net architecture and its variants, have been extensively deployed. As computational resources have expanded, the field has seen a rapid influx of highly complex architectures, including 3D convolutions, Attention U-Nets, and recently, Vision Transformers (ViTs) such as UNETR and Swin-UNet. While these models excel at capturing long-range global dependencies, they impose severe memory constraints. Consequently, a common practice in medical imaging pipelines is to aggressively downsample the input volumes (e.g., from $512 \times 512$ native resolution to $256 \times 256$ or $128 \times 128$) to fit entire 3D volumes or complex multi-layer transformer blocks into GPU memory.

In this paper, we challenge this paradigm. We hypothesize that for small, low-contrast anatomical structures like the pancreas, the loss of fine-grained local texture and boundary sharpness caused by downsampling outweighs the benefits of advanced architectural complexity or global context.

\subsection{Contributions}
\begin{itemize}
    \item \textbf{Resolution Ablation Study:} We provide rigorous empirical evidence demonstrating that downsampling CT slices causes a collapse in segmentation performance for the pancreas, proving that spatial resolution is the primary performance driver.
    \item \textbf{Domain-Specific Optimization:} We quantify the impact of optimized HU windowing ([-125, 275]), showing a +4.9\% absolute Dice gain over standard windowing.
    \item \textbf{SOTA vs. Complexity:} We demonstrate that a simple, patch-based CNN operating at native resolution outperforms complex Vision Transformers forced to operate under standard memory-constrained conditions.
    \item \textbf{Annotation Efficiency via SSL:} We integrate our resolution-preserving framework with advanced Semi-Supervised Learning (SSL) methods (Mean Teacher, Uncertainty-Aware Mean Teacher, and Cross-Pseudo Supervision). We demonstrate that combining high spatial resolution with consistency regularization recovers near-fully-supervised performance while using only a fraction of labeled data.
\end{itemize}

\section{Methodology}

\subsection{Data and Preprocessing}
We utilized the public NIH Pancreas-CT dataset (Task07\_Pancreas from the Medical Segmentation Decathlon), consisting of 82 contrast-enhanced abdominal CT volumes. 

\subsubsection{Domain-Specific HU Windowing}
Standard abdominal windowing typically clips HU values to $[-100, 240]$. However, the pancreas often exhibits subtle intensity gradients that blend into surrounding fat and vascular structures. We empirically identified and applied an optimized window of $[-125, 275]$. The values were subsequently min-max normalized to $[0, 1]$. 

\subsubsection{Resolution Preservation Strategy}
Instead of downsampling the native $512 \times 512$ axial slices, we employed a patch-based extraction strategy. During training, $256 \times 256$ patches were randomly cropped from the native resolution slices. This allowed the network to learn from uncompromised pixel data while maintaining manageable memory footprints. During inference, a sliding-window approach with a stride of 128 was used to reconstruct the full $512 \times 512$ probability map, utilizing Gaussian weighting to smooth overlapping predictions.

\subsection{Architectural Baselines}
To isolate the effect of resolution versus architecture, we evaluated the following configurations:
\begin{enumerate}
    \item \textbf{Patch-based SOTA (Our Method):} A standard 4-level U-Net trained on native-resolution $256 \times 256$ patches.
    \item \textbf{Ablation 256x256:} A standard U-Net trained on full CT slices downsampled to $256 \times 256$.
    \item \textbf{Ablation 128x128:} A standard U-Net trained on full CT slices downsampled to $128 \times 128$.
    \item \textbf{Transformer Baseline:} A Vision Transformer U-Net (TransUNet style) \cite{transunet} trained on the native patches to evaluate the necessity of complex self-attention mechanisms.
    \item \textbf{nnU-Net Baseline:} The established medical image segmentation state-of-the-art framework \cite{nnunet}, utilizing a self-configuring 3D full-resolution architecture. We trained the `3d_fullres` configuration on our dataset to provide an aggressive fully supervised benchmark.
\end{enumerate}

\subsection{Semi-Supervised Learning (SSL) Framework}
To address the clinical bottleneck of voxel-level annotation, we investigated the "break-even point" of data efficiency using three state-of-the-art SSL paradigms:
\begin{enumerate}
    \item \textbf{Mean Teacher (MT):} A consistency-regularization approach \cite{tarvainen2017mean} where a student model learns from the exponential moving average (EMA) of its own weights (the teacher). We enforce Mean Squared Error (MSE) consistency between their predictions on unlabeled data.
    \item \textbf{Uncertainty-Aware Mean Teacher (UA-MT):} An extension of MT \cite{uamt} that utilizes Monte Carlo (MC) Dropout to estimate the teacher's epistemic uncertainty. The student is forced to learn only from regions where the teacher is highly confident, which is crucial for the ambiguous boundaries of the pancreas.
    \item \textbf{Cross-Pseudo Supervision (CPS):} A mutual-learning approach \cite{cps} where two identically structured networks with different initializations generate pseudo-labels to supervise each other.
\end{enumerate}
We evaluated these methods using varying labeled data ratios (10\%, 25\%, and 50\%) against a fully supervised upper bound.

\section{Results and Discussion}

\subsection{The Impact of Domain-Specific Windowing}
Our initial experiments isolated the effect of HU windowing using our optimized patch-based model. Evaluating on a held-out validation set, the standard $[-100, 240]$ window achieved an average Dice coefficient of $0.7653$. Shifting to the pancreas-specific $[-125, 275]$ window improved the average Dice to $0.8147$. This absolute gain of nearly 5\% underscores the sensitivity of CNNs to raw input distributions in medical imaging.

\subsection{Resolution Ablation: The Collapse of the Pancreas}
We conducted a rigorous ablation study to quantify the cost of downsampling. 
The $128 \times 128$ full-slice model converged rapidly but stagnated at a Validation Intersection-over-Union (IoU) of $0.3036$. The $256 \times 256$ full-slice model showed only a marginal improvement, reaching a Validation IoU of $0.3113$. 

Crucially, both downsampled models dramatically underperformed our patch-based SOTA model (0.849 Dice, $\sim 0.69$ IoU equivalent). This massive performance cliff indicates that the pancreas is effectively "lost" to the network when native resolution is compromised. The global context provided by seeing the entire abdomen in the downsampled slices could not compensate for the loss of local boundary texture.

\subsection{Architectural Complexity: 3D Models and ViTs vs. High-Res 2D CNNs}
We extended our architectural comparison to evaluate the industry-standard nnU-Net framework, specifically its \textit{3d\_fullres} configuration which natively trains on 3D volumetric patches. After exhaustive convergence (1000 epochs), the 3D nnU-Net achieved an average Pancreas Dice of 0.8225. While this is an excellent performance, it still fell slightly short of our optimized 2D Patch-based CNN (0.849 Dice). This validates our hypothesis: maintaining high-resolution 2D spatial features and utilizing an aggressive overlapping sliding window during inference is highly competitive with, and often superior to, learning native 3D convolutions for small, low-contrast targets.

Furthermore, we evaluated a Vision Transformer (ViT) architecture (TransUNet style) trained on the same native-resolution 2D patches as our SOTA CNN baseline. The Transformer baseline reached a Validation IoU of 0.390. Despite its complex self-attention mechanisms and significantly higher computational cost, it drastically underperformed the simpler patch-based CNN (which achieved an equivalent IoU of $\sim 0.69$). This confirms that for highly variable organs like the pancreas, learning robust local textures via convolutions is vastly more effective than attempting to capture long-range global dependencies without massive pre-training.

\subsection{Annotation Efficiency: The SSL Break-Even Point}

To evaluate annotation efficiency, we trained Semi-Supervised Learning (SSL) methods—Mean Teacher (MT), Cross-Pseudo Supervision (CPS), and Uncertainty-Aware Mean Teacher (UA-MT)—using 10\%, 25\%, and 50\% of the available labeled data. The validation Intersection over Union (IoU) results are presented in Table~\ref{tab:annotation_efficiency}.

\begin{table}[h]
\centering
\caption{Annotation Efficiency (Validation IoU)}
\label{tab:annotation_efficiency}
\begin{tabular}{lccc}
\hline
\textbf{Method} & \textbf{10\% Labeled} & \textbf{25\% Labeled} & \textbf{50\% Labeled} \\
\hline
Mean Teacher    & 0.1939               & 0.3956               & 0.4571               \\
CPS             & 0.2456               & 0.3328               & Pending*             \\
UA-MT           & \textbf{0.3965}      & \textbf{0.4836}      & Pending*             \\
\hline
\multicolumn{4}{l}{\footnotesize *Training is currently in progress.}
\end{tabular}
\end{table}

These results confirm that UA-MT consistently dominates other SSL approaches, particularly in extremely low-data regimes. At just 10\% labeled data, UA-MT approaches the performance of Mean Teacher at 25\%, effectively halving the required annotation effort. This substantial gain is driven by its ability to rigorously filter ambiguous boundary regions during pseudo-labeling.

\subsection{Volumetric Performance and Clinical Relevance}
To validate the clinical utility of our SSL framework, we performed full 3D volumetric inference using a sliding-window approach on the champion models. As shown in Table~\ref{tab:volumetric_dice}, our Fully Supervised SOTA achieved a Mean 3D Dice of 0.849. Remarkably, the UA-MT model trained on only 25\% of the labels reached a 3D Dice of 0.7241, while the Mean Teacher model at 50\% labels reached 0.7585. 

\begin{table}[h]
\centering
\caption{Volumetric Performance (3D Test Dice) for Champion Models}
\label{tab:volumetric_dice}
\begin{tabular}{lc}
\hline
\textbf{Configuration} & \textbf{Mean 3D Dice} \\
\hline
Fully Supervised (100\%) & 0.8490 \\
Mean Teacher (50\%)      & 0.7585 \\
UA-MT (25\%)            & 0.7241 \\
\hline
\end{tabular}
\end{table}

The results demonstrate that by preserving native resolution, even semi-supervised models can capture the fine-grained morphology of the pancreas. Notably, on "stable" cases (e.g., Pancreas\_001), the 50\% Mean Teacher model achieves a Dice of 0.9002, nearly identical to the supervised upper bound (0.9018), suggesting that for a large subset of clinical cases, 50\% annotation is sufficient for high-fidelity segmentation.

\section{Conclusion}
[To be completed as Phase 1 finalizes.]

\bibliographystyle{IEEEtran}
\bibliography{references}

\end{document}
